% -----------------------------------------------------------
\section{Order (noun)}
% -----------------------------------------------------------
	\label{ORDER}
	
% % ----------------------
% \subsection{See also}
% % ----------------------

% ----------------------
\subsection{1828 American Dictionary of the English Language} % *1828
% ----------------------
	\label{ORDER-1828}
	
	% -----
	\subsubsection{Noun}
	% -----

		\begin{compactitem}
			\item[1] Regular disposition or methodical arrangement of things; a word of extensive application; as the order of troops or parade; the order of books in a library; the order of proceedings in a legislative assembly. order is the life of business.
			\item[2] Proper state; as the muskets are all in good order When the bodily organs are in order a person is in health; when they are out of order he is indisposed.
			\item[3] Adherence to the point in discussion, according to established rules of debate; as, the member is not in order that is, he wanders from the question.
			\item[4] Established mode of proceeding. The motion is not in order.
			\item[5] Regularity; settled mode of operation. This fact could not occur in the order of nature; it is against the natural order of things.
			\item[6] Mandate; precept; command; authoritative direction. I have received an order from the commander in chief. The general gave orders to march. There is an order of council to issue letters of marque.
			\item[7] Rule; regulation; as the rules and orders of a legislative house.
			\item[8] Regular government or discipline. It is necessary for society that good order should be observed. The meeting was turbulent; it was impossible to keep order
			\item[9] Rank; class; division of men; as the order of nobles; the order of priests; the higher orders of society; men of the lowest order; order of knights; military orders, etc.
			\item[10] A religious fraternity; as the order of Benedictines.
			\item[11] A division of natural objects, generally intermediate between class and genus. The classes, in the Linnean artificial system, are divided into orders, which include one or more genera. Linne also arranged vegetables, in his natural system, into groups of genera, called order In the natural system of Jussieu, orders are subdivisions of classes.
			\item[12] Measures; care. Take some order for the safety and support of the soldiers.
			\item[13] In rhetoric, the placing of words and members in a sentence in such a manner as to contribute to force and beauty of expression, or to the clear illustration of the subject.
			\item[14] The title of certain ancient books containing the divine office and manner of its performance.
			\item[15] In architecture, a system of several members, ornaments and proportions of columns and pilasters; or a regular arrangement of the projecting parts of a building, especially of the columns, so as to form one beautiful whole. 
		\end{compactitem}

% % ----------------------
% \subsection{Oxford English Dictionary} % *OED
% % ----------------------
	% \label{ORDER-OED}

	% \begin{compactitem}
		% \item[]
	% \end{compactitem}

% ----------------------
\subsection{Scriptural Usage} % *SCRIPTURES
% ----------------------
	\label{ORDER-SCRIPTURES}

	% % -----
	% \subsubsection{Old Testament} % *OT
	% % -----
		% \label{ORDER-OT}
	
		% % --
		% \paragraph{KJV}
		% % --
			% \label{ORDER-OT-KJV}
	
			% \begin{tabular}{ll}
				% Total occurrences in the OT & 
			% \end{tabular}
			
			% \begin{compactdesc}
				% \item[]
			% \end{compactdesc}
			
		% % --
		% \paragraph{Hebrew Bible}
		% % --
			% \label{ORDER-HEBREW}
		
			% %
			% \subparagraph{\texorpdfstring{\he{sample word}}{}} % paste actual hebrew text here
			% %
				% \label{PAGA} % whatever is in step bible without periods
			
				% \begin{compactdesc}
					% \item[HALOT , PDF ] \textbf{} % Use the 1994 PDF
						% \begin{compactitem}
							% \item[1]
						% \end{compactitem}
					% \item[BDB , PDF ] \textbf{}
						% \begin{compactitem}
							% \item[1]
						% \end{compactitem}
					% \item[DCH PDF ] \textbf{} % don't include the actual page number, they repeat from volume to volume
						% \begin{compactitem}
							% \item[1]
						% \end{compactitem}
				% \end{compactdesc}
				
				% \begin{compactdesc}
					% \item[Some OT uses] \textbf{}
						% \begin{compactdesc}
							% \item[]
						% \end{compactdesc}
				% \end{compactdesc}
			
		% % --
		% \paragraph{Septuagint}
		% % --
			% \label{ORDER-LXX}
		
			% %
			% \subparagraph{\texorpdfstring{\gr{sample word}}{}}
			% %
		
	% -----
	\subsubsection{New Testament} % *NT
	% -----
		\label{ORDER-NT}
	
		% % --
		% \paragraph{KJV}
		% % --
			% \label{ORDER-NT-KJV}		
	
			% \begin{tabular}{ll}
				% Total occurrences in the NT & 
			% \end{tabular}
			
			% \begin{compactdesc}
				% \item[]
			% \end{compactdesc}
			
		% --
		\paragraph{Greek New Testament} % *GREEK
		% --
			\label{ORDER-GREEK}
		
			%
			\subparagraph{\texorpdfstring{\gr{τάξις}}{}}
			%
				\label{TAXIS}
				
				\begin{tabular}{ll}
					Total occurrences in the NT & 9 % ECGNT 930b
				\end{tabular}
				
				\begin{compactdesc}
					\item[All NT uses] \textbf{} % Change to "all" if they're all listed
						\begin{compactdesc}
							\item[{\hyperref[LUKE-1-8]{Luke 1:8}}]
							\item[{\hyperref[1CORINTHIANS-14-40]{1 Corinthians 14:40}}]
							\item[{\hyperref[COLOSSIANS-2-5]{Colossians 2:5}}]
							\item[{\hyperref[HEBREWS-5-6]{Hebrews 5:6}}]
							\item[{\hyperref[HEBREWS-5-10]{Hebrews 5:10}}]
							\item[{\hyperref[HEBREWS-6-20]{Hebrews 6:20}}]
							\item[{\hyperref[HEBREWS-7-11]{Hebrews 7:11}}] \textbf{(2)}
							\item[{\hyperref[HEBREWS-7-17]{Hebrews 7:17}}]
						\end{compactdesc}
				\end{compactdesc}
			
				\begin{compactdesc}
				
					\item[BDAG 879a] \phantom{.}
						\begin{compactitem}
							\item[1] \textbf{an arrangement of things in sequence, \textit{fixed succession/order}}
							\item[2] \textbf{a state of bood order, \textit{order, proper procedure}}
							\item[3] \textbf{an assigned station or rank, \textit{position, post}}
							\item[4] \textbf{an arrangement in which someone or something functions, \textit{arrangement, nature, manner, condition, outward aspect}}...Perhaps it is in this way that Hebrews understood \hyperref[PSALMS-110-4]{Ps 109:4b}, which the author interprets to mean that Jesus was a high priest \gr{κατὰ τὴν τάξιν Μελχισέδεκ} \textit{according to the nature of} = \textit{just like Melchizedek} i.e. like the type of arrangement made for the functioning of Melchizedek: \hyperref[HEBREWS-5-6]{5:6}, \hyperref[HEBREWS-5-10]{10}; \hyperref[HEBREWS-6-20]{6:20}; \hyperref[HEBREWS-7-11]{7:11a}, \hyperref[HEBREWS-7-17]{17}, \hyperref[HEBREWS-7-21]{21} \textit{varia lectio}. In any case the reference is not only to the higher `rank', but also to the entirely different nature of Melchizedek's priesthood as compared with that of Aaron \hyperref[HEBREWS-7-11]{7:11b}.
						\end{compactitem}
						
					\item[Brill 2082b, PDF 2146] \phantom{.}
						\begin{compactitem}
							\item[A] \textit{military} \textbf{arrangement, formation, disposition} \textit{of an army}; \textbf{battle order}; \textbf{row} \textit{or} \textbf{line} of soldiers; \textbf{group, crowd}
							\item[B] \textit{usually} \textbf{disposition}; \textbf{arrangement} \textit{of the elements of a speech}; \textbf{orderliness, regularity}; \textbf{order, prescription, provision}; \textbf{rule}; \textit{religious} \textbf{order}
								\begin{compactitem}
									\item The Hebrews usage is mentioned here.
								\end{compactitem}
							\item[C] \textit{figurative} \textbf{place, position, state, part}; \textbf{order, social class}; \textbf{political system, constitution}
							\item[D] \textit{medical} \textbf{reduction}
							\item[E] \textit{plural} \textbf{acts, documents}
						\end{compactitem}
						
					\item[CGL 1358a, PDF 1398] \phantom{.}
						\begin{compactitem}
							\item[1] (military) process of arranging in order, \textbf{organising, disposition}
							\item[2] arranged order, \textbf{formation, battle order}
							\item[3] (often plural) \textbf{rank, line} (of troops)
							\item[4] \textbf{contingent} (of troops)
							\item[5] \textbf{station, post} (of troops)
							\item[6] (in non-military contexts) \textbf{order, arrangement} (of one's personal life, things in the physical world); \textbf{organisation, system} (of government); (as a more general organising principle, statements associated with law) \textbf{good order, regularity}
							\item[7] (law) imposed arrangement, \textbf{ordinance, regulation}
							\item[8] \textbf{designated location} or \textbf{spatial position}
							\item[9] (in political context) \textbf{chosen position, stance}
							\item[10] \textbf{class, category} (of persons, referring to their character or social position)
						\end{compactitem}
						
					\item[LSJ 1756b, PDF 1827] \phantom{.}
						\begin{compactitem}
							\item[I] in military senses
							\item[I.5] \textit{post} or \textit{place in the line} of battle
							\item[II.1] generally, \textit{arrangement, order}
							\item[II.2] \textit{order, regularity}
							\item[II.3] \textit{ordinance}
							\item[III] metaphorical from I.5, \textit{post, rank, position, station}
							\item[IV.1] \textit{order, class} of men
						\end{compactitem}
						
					\item[EDNT 3:333, PDF 1447] \phantom{.}
						\begin{compactitem}
							\item[2] It refers to \textit{order} within the congregation...Luke 1:8 refers to the \textit{sequence/turn} of service of the priestly \gr{ἐφημερία}.
							\item[3] \gr{τάξις} refers also to a particular \textit{condition, nature, manner}...following Ps 109:4 LXX, `according to the \textit{nature} of Melchizedek,' Heb 5:6, 10; 6:20; 7:11c, 17 is interpreted by \gr{κατὰ τὴν ὁμοιότητα Μελχισέδεκ} in 7:15, in contrast to \gr{κατὰ τὴν τάξιν Ἀαρὼν} in 7:11d... W. R. G. Loader...interprets the alternative \textit{ways} of functioning as a priest, established either through \gr{νόμος} or through \gr{ὁρκωμοσία}, in the sense of an old and new \textit{order}, qualified by mortality on the one hand, and indestructible life on the other.
							\item See this entry \href{https://drive.google.com/file/d/1xdzDS1V5SS8b_kukGSVS9kZfrgnQmfwt/view?usp=sharing}{here}.
						\end{compactitem}
						
					\item[TDNT] \phantom{.}
						\begin{compactitem}
							\item DNA
								\begin{compactitem}
									\item See 8:27, note 1---is it possible they thought the word didn't appear in the NT? Hard to believe.
								\end{compactitem}
						\end{compactitem}
						
					% \item[TDNTA] \phantom{.}
						% \begin{compactitem}
							% \item 
						% \end{compactitem}
						
					% \item[NIDNTT] \phantom{.}
						% \begin{compactitem}
							% \item 
						% \end{compactitem}
						
					\item[NIDNTTE] \phantom{.}
						\begin{compactitem}
							\item DNA
						\end{compactitem}
						
				\end{compactdesc}
		
	% % -----
	% \subsubsection{Book of Mormon} % *BOM
	% % -----
		% \label{ORDER-BOM}
	
		% \begin{tabular}{ll}
			% Total occurrences in the BOM & 
		% \end{tabular}
		
		% \begin{compactdesc}
			% \item[]
		% \end{compactdesc}
		
	% -----
	\subsubsection{Doctrine and Covenants} % *DC
	% -----
		\label{ORDER-DC}
	
		% \begin{tabular}{ll}
			% Total occurrences in the DC & 
		% \end{tabular}
		
		\begin{compactdesc}
			\item[{\hyperref[DC94-6]{DC 94:6}}] Need to look at images of the actual building to understand this scripture better.
		\end{compactdesc}
		
	% % -----
	% \subsubsection{Pearl of Great Price} % *PGP
	% % -----
		% \label{ORDER-PGP}
	
		% \begin{tabular}{ll}
			% Total occurrences in the PGP & 
		% \end{tabular}
		
		% \begin{compactdesc}
			% \item[]
		% \end{compactdesc}
		
% ----------------------
\subsection{Key Phrases} % *KEYPHRASES *PHRASES
% ----------------------
	\label{ORDER-PHRASES}

	\begin{compactdesc}
		\item[holy order] \textbf{\#times} | \hyperref[2NEPHI-6-2]{2 Nephi 6:2}
			% \begin{compactdesc}
				% \item[Bible anchor verses] \phantom{.}
					% \begin{compactdesc}
						% \item[Romans 5:5] \gr{}
					% \end{compactdesc}
				% \item[Commentary] ---
			% \end{compactdesc}
	\end{compactdesc}
	
% % ----------------------
% \subsection{Bible Dictionaries} % *DICTIONARIES
% % ----------------------
	% \label{ORDER-BD}

% % ----------------------
% \subsection{Restoration Usage} % *RESTORATION
% % ----------------------
	% \label{ORDER-RESTORATION}

	% % -----
	% \subsubsection{Joseph Smith} % *JS
	% % -----
		% \label{ORDER-JS}
	
		% \begin{compactdesc}
			% \item[{\hyperref[KEY]{Date/Source}}]
		% \end{compactdesc}
		
	% % -----
	% \subsubsection{Temple Ordinances}
	% % -----
		% \label{ORDER-TEMPLE}
	
		% \begin{compactdesc}
			% \item[{\hyperref[KEY]{Date/Source}}]
		% \end{compactdesc}
	
	% % -----
	% \subsubsection{Brigham Young} % *BY
	% % -----
		% \label{ORDER-BY}
	
		% \begin{compactdesc}
			% \item[{\hyperref[KEY]{Date/Source}}]
		% \end{compactdesc}
	
% % ----------------------
% \subsection{Commentary and Studies} % *COMMENTARY *STUDIES
% % ----------------------
	% \label{ORDER-COMMENTARY}

	% % -----
	% \subsubsection{Key Points}
	% % -----
		% \label{ORDER-KEYS}
	
		% \begin{compactitem}
			% \item
		% \end{compactitem}
		
	% % -----
	% \subsubsection{Sample Study}
	% % -----
		% \label{ORDER-study}